
\chapter{Investigating Indoor Experience Through Affective Interaction}
\label{ch:affective-interaction}

\begin{synopsis}
\synopsislead{How can Affective Interaction help us investigate}{the complexity of lived experience in interactive indoor environments? The preceding review identified emotional experience as a substantial strand of research, while showing a need for methods that can retain the situated complexity through which indoor conditions come to matter. This chapter examines Affective Interaction (AfI), developed within HCI, as a methodological lens for addressing this need.}

Smart building research has commonly investigated occupant experience through measurable dimensions of environmental comfort, including thermal, visual, acoustic, and respiratory conditions. Architectural accounts of inhabitation, however, draw attention to experience as bodily, situated, interpretive, and socially embedded. AfI provides a way of bringing these concerns into HBI by treating affect as part of how people perceive, interpret, and establish relations with their surroundings.

The chapter investigates this proposition through two exploratory studies in a smart university building. An in-situ survey examines what emotion- and comfort-framed prompts draw into accounts of the same indoor experience, while a semi-structured building walk follows how occupants interpret and adjust to their surroundings across space and time. The analysis identifies five \emph{Ways} in which occupants experience and interpret smart environments, from reading bodily cues and anticipating disruption to constructing spatial purpose, adjusting to changing conditions, and making sense of building automation. From these findings, the chapter develops six \emph{Affective Practices} and six corresponding \emph{Interaction Qualities}.

Within the trajectory of the dissertation, this chapter addresses \rqref{affective} by showing what an affective lens can reveal about indoor experience and how this knowledge can inform HBI research and design. AfI directs attention to the relations among bodily sensation, environmental conditions, activity, social situation, temporality, and interpretations of building behaviour through which experience acquires significance. It therefore establishes the methodological basis for the dissertation's affective orientation before the following chapters turn from investigating lived experience to developing experiential possibilities for future interactive indoor environments.
\end{synopsis}

\chaptersource{%
This chapter is based on: Shruti Rao, Judith Good, and Hamed Alavi.
\emph{Emotion in Smart Buildings: Can Affective Interaction Shape the Smart Agenda in Architecture?}
Presented at CHI 2026, Barcelona, Spain. Honourable Mention.%
}

\clearpage

\import{papers/02-afi/}{thesis-body.tex}
\FloatBarrier
\clearpage
\begin{chapterclosing}{What this chapter establishes}

This chapter establishes Affective Interaction as a methodological lens for investigating lived experience in interactive indoor environments. Across the two studies, affect directed attention beyond environmental appraisal towards how occupants sense, interpret, anticipate, and adjust to their surroundings in relation to bodily state, activity, spatial conditions, other people, time, and building behaviour.

The findings show that these relations are central to how indoor conditions acquire significance. The five \emph{Ways} describe how occupants interpret and act within smart environments, while the six \emph{Affective Practices} articulate recurrent processes through which they orient, select, coordinate, adjust, make sense of automation, and settle within their surroundings. The corresponding \emph{Interaction Qualities} carry this experiential knowledge into design by identifying qualities through which interactive environments might support these practices. 

Within the dissertation, this chapter therefore develops the methodological basis of affective HBI. AfI provides a way to investigate experience without reducing it to isolated environmental variables or momentary affective states, while retaining the bodily, spatial, social, temporal, and interpretive relations through which occupants experience indoor environments. The chapter also establishes a connection between investigating experience and designing for it: situated accounts of how people relate to existing environments can inform the qualities considered in future interactive environments. 

\chaptercarryforward{The following chapter develops this design direction through multisensory biophilic design, drawing on architectural expertise to examine what forms of experience future interactive indoor environments might support.}
\end{chapterclosing}
% \begin{chapterclosing}{What this chapter establishes}
% This chapter demonstrates how Affective Interaction can be used to investigate the complexity of lived indoor experience. The emotion-framed survey drew attention to situated affect, bodily sensation, and social context, while the building walk traced how these experiences developed through bodily cues, anticipation, spatial interpretation, adjustment, and sense-making around building automation.

% Across these studies, affect provides a means of examining how indoor conditions acquire significance within particular situations. The resulting Affective Practices articulate recurrent ways in which occupants orient to, negotiate, and settle within their surroundings, while the corresponding Interaction Qualities translate these processes into considerations for interactive systems. Together, they extend the account of human experience developed in the preceding chapter from a field-level specification towards a situated account of how experience is sensed, interpreted, and acted upon.

% Within the dissertation, this establishes the methodological component of affective HBI: attending to the relations through which indoor experience becomes meaningful and using these relations to inform design. The next stage develops the design component by asking what kinds of experiential possibilities future interactive indoor environments might support.

% \chaptercarryforward{The following chapter turns from investigating existing indoor experience to designing future interactive environments, drawing on architectural expertise to examine the possibilities of multisensory biophilic design.}
% \end{chapterclosing}